% from ega3-1-fr.tex
\oldpage[0\textsubscript{III}]{37}
propriétés suivantes~:
\begin{enumerate}
  \item[$1^\circ$] Tout $Q^i$ est un objet injectif de $\mathcal C$~;
  \item[$2^\circ$] Pour tout $i\geq0$, on a
    $\mathrm Z^i(Q^\bullet)=\mathrm B^i(Q^\bullet)$, et
    $\mathrm Z^i(Q^\bullet)$ est facteur direct de $Q^i$.
\end{enumerate}
Alors~:
\begin{enumerate}
  \item[(i)] Tout $Q^\bullet\in\mathfrak J$ est un objet injectif de
    $\mathcal K$.
  \item[(ii)] Tout objet de $\mathcal K$ est isomorphe à un sous-complexe
    d'un complexe appartenant à $\mathfrak J$.
\end{enumerate}
\end{lemma}

\begin{proof}
\textup{(i)} Soient $A^\bullet=(A^i)$ un objet de $\mathcal K$,
$A'^\bullet=(A'^i)$ un sous-objet de $A^\bullet$,
$Q^\bullet=(Q^i)$ un objet de $\mathfrak J$, et supposons donné un morphisme
$f=(f^i):A'^\bullet\to Q^\bullet$, qu'il s'agit de prolonger en un morphisme
$g=(g^i):A^\bullet\to Q^\bullet$. Nous utiliserons le langage de la catégorie
des modules pour simplifier (cf. [27]).

Identifions $Q^i$ à
$\mathrm B^i(Q^\bullet)\oplus\mathrm B^{i+1}(Q^\bullet)$~; procédons par
récurrence sur $i$, et supposons donc les $g^j$ définis pour $j<i$,
compatibles avec les opérateurs de dérivation $d^j:A^j\to A^{j+1}$ et
$d'^j:Q^j\to Q^{j+1}$ pour $j<i-1$ et tels en outre que~:
\begin{enumerate}
  \item[$1^\circ$] $g^{i-1}(\mathrm Z^{i-1}(A^\bullet))\subset
    \mathrm Z^{i-1}(Q^\bullet)$~;
  \item[$2^\circ$] Si on pose $C^j=(d^j)^{-1}(A'^{j+1})$ pour tout $j$, alors
    $d'^{i-1}\circ g^{i-1}$ coïncide avec $f^i\circ d^{i-1}$ sur $C^{i-1}$.
\end{enumerate}
Le morphisme $f^i:A'^i\to Q^i$ donne par composition avec les projections
deux morphismes
$f'^i:A'^i\to\mathrm B^i(Q^\bullet)$ et
$f''^i:A'^i\to\mathrm B^{i+1}(Q^\bullet)$. Comme
$d'^{i-1}\circ g^{i-1}$ applique $A^{i-1}$ dans
$\mathrm B^i(Q^\bullet)$ et s'annule dans $\mathrm Z^{i-1}(A^\bullet)$, il
définit un morphisme
$h^i:\mathrm B^i(A^\bullet)\to\mathrm B^i(Q^\bullet)$, et puisque
$d'^{i-1}\circ g^{i-1}$ coïncide avec $f^i\circ d^{i-1}$ sur $C^{i-1}$,
$h^i$ coïncide avec $f'^i$ dans
$\mathrm B^i(A^\bullet)\cap A'^i$. Comme $\mathrm B^i(Q^\bullet)$, facteur
direct de $Q^i$, est injectif, il y a un morphisme
$g'^i:A^i\to\mathrm B^i(Q^\bullet)$ qui coïncide avec $h^i$ dans
$\mathrm B^i(A^\bullet)$ et avec $f'^i$ dans $A'^i$. Considérons d'autre part
le morphisme
$f''^{i+1}\circ d^i:C^i\to\mathrm B^{i+1}(Q^\bullet)$, qui s'annule dans
$\mathrm Z^i(A^\bullet)$~; comme $\mathrm B^{i+1}(Q^\bullet)$ est injectif,
il y a un morphisme $g''^i:A^i\to\mathrm B^{i+1}(Q^\bullet)$, qui coïncide
avec $f''^{i+1}\circ d^i$ dans $C^i$ et avec $0$ dans
$\mathrm Z^i(A^\bullet)$. Il suffit alors de prendre $g^i=g'^i+g''^i$ pour
pouvoir poursuivre la récurrence.

\textup{(ii)} Pour plonger $A^\bullet=(A^i)$ dans un complexe appartenant à
$\mathfrak J$, on prend pour chaque $i\geq1$ un objet injectif $Q'^i$ de
$\mathcal C$ tel qu'il existe une injection $f'^i:A^i\to Q'^i$. Posons alors
\[
  Q^i=0\quad(i<0),\qquad Q^0=Q'^1,\qquad
  Q^i=Q'^i\oplus Q'^{i+1}\quad(i\geq1),
\]
avec l'opérateur de dérivation évident. Posons
$f^i=f'^i+(f'^{i+1}\circ d^i)$ pour tout $i$ (avec $f'^i=0$ pour $i\leq0$)~;
il est immédiat que $f^i$ est injectif pour tout $i$, et que les $f^i$ sont
compatibles avec les opérateurs de dérivation.
\end{proof}

\begin{corollary}[11.5.2.2]
\label{0.11.5.2.2-fr}
Tout objet $K^\bullet$ de $\mathcal K$ admet une résolution droite formée
d'objets de $\mathfrak J$. Si $L^{\bullet\bullet}$ est une telle résolution,
pour toute résolution $L'^{\bullet\bullet}$ de $K^\bullet$, formée d'objets
de $\mathcal K$, il y a un morphisme de bicomplexes
$F:L'^{\bullet\bullet}\to L^{\bullet\bullet}$ compatible avec les
augmentations, et deux tels morphismes $F$, $F'$ sont homotopes.
\end{corollary}

\begin{proof}
Ce n'est autre que (M, V, 1.1 a)) appliqué à la catégorie abélienne
$\mathcal K$.
\end{proof}

\begin{env}[11.5.2.3]
\label{0.11.5.2.3-fr}
Ces préliminaires étant posés, considérons une résolution de
Cartan-Eilenberg injective $L'^{\bullet\bullet}$ de $K^\bullet$ et une
résolution $L^{\bullet\bullet}$ de $K^\bullet$ formée d'objets de
$\mathfrak J$, et montrons que l'on a un isomorphisme
\[
  \mathrm H^\bullet(T(L'^{\bullet\bullet}))\xrightarrow{\sim}
  \mathrm H^\bullet(T(L^{\bullet\bullet})).
\]
On déduit en effet de (\hyperref[0.11.5.2.2-fr]{11.5.2.2}) un morphisme de
bicomplexes $T(L'^{\bullet\bullet})\to T(L^{\bullet\bullet})$, et par suite
un morphisme
$'E(T(L'^{\bullet\bullet}))\to{}'E(T(L^{\bullet\bullet}))$ des premières
suites spectrales de ces bicomplexes. Comme en vertu de
(\hyperref[0.11.3.3-fr]{11.3.3}), ces suites spectrales sont régulières, il
suffit (\hyperref[0.11.1.5-fr]{11.1.5}) de voir que le morphisme précédent est
un isomorphisme pour les termes $E_2$, ou encore que
$\mathrm H_{\mathrm{II}}^q(T(L^{i,\bullet}))$ est égal à
$\mathrm R^qT(K^i)$~; comme $L^{i,\bullet}$ est une résolution droite de
$K^i$, on est ramené à prouver le
\end{env}

